%-------------------------------------------------------------------------------
%	EXPERIENCE
%-------------------------------------------------------------------------------
\cvsection{Experience}

\begin{cventries}

\cventry
    {Engineering Manager: Data \& Bioinformatics}
    {Debiopharm}
    {Lausanne, Switzerland}
    {July 2025 - Present}
    {
      \begin{cvitems}
        \item Lead a cross-functional team of data engineers, bioinformaticians \& data scientists; set up a bioinformatics \& data science service model for the Translational Medicine \& Antibody Engineering teams, with a service catalogue, sprint-based intake \& quarterly prioritisation
        \item Own preclinical analytics across in-vitro assays (cytotoxicity, internalization, binding) \& in-vivo efficacy studies, delivering a dozen+ interactive dashboards used daily by translational scientists
        \item Lead transcriptomics biomarker discovery \& validation for clinical-stage oncology programs: RNA-seq QC, batch correction, differential expression, pathway enrichment \& ML gene signatures
        \item Deliver in-silico antibody developability \& CMC risk assessment for antibody-drug conjugate (ADC) programs using protein language models, structure prediction \& sequence-liability scanning
        \item Led the selection of \& migration to a new NGS analysis platform, bringing omics processing in-house; built RNA-seq, variant-calling \& ctDNA pipelines with automated, secure transfers from CROs \& sequencing vendors
        \item Defined the preclinical data-flow \& governance framework across the data repository, ELN \& metadata systems, aligning Data Management, IT, Translational Medicine, Quality \& wet-lab teams
        \item Initiated applied AI projects: generative models \& Monte Carlo simulation for bioequivalence trial design, \& a digital-pathology pipeline for IHC Ki-67 scoring
      \end{cvitems}
    }

\cventry
    {Data Engineering Team Lead}
    {}
    {}
    {July 2022 - June 2025}
    {
      \begin{cvitems}
        \item Architected \& led, from scratch, the Central Data Repository: the company's hub for all clinical \& non-clinical data, used by teams across several functions \& all R\&D programs
        \item Built the scientific analytics layer of the data fabric (SAS Viya on Azure Kubernetes, Posit suite, Terraform IaC), serving 200+ users \& 100+ data products under GxP validation
        \item Delivered DropZone, a secure data-exchange platform for partners \& vendors, in five months; used for in- \& out-licensing due diligence
        \item Productionised privacy-aware LLM applications for document analysis \& protocol writing, \& an entity-recognition platform for scientific literature monitoring
        \item Automated GxP computerised system validation, infrastructure qualification \& security audits through DevOps; led technical due diligence of healthcare start-ups for the corporate venture fund
      \end{cvitems}
    }

\cventry
    {Senior Data Engineer}
    {SOPHiA GENETICS}
    {Lausanne, Switzerland}
    {April 2021 - June 2022}
    {
      \begin{cvitems}
        \item Built terabyte-scale ETL pipelines turning raw genomics data into analytics-ready formats, with APIs \& CLI tools that opened the data to internal teams
        \item Introduced regression testing for variant-detection pipelines \& TLA+ formal specifications to verify the correctness of ETL operations
      \end{cvitems}
    }

\cventry
    {Software Engineer | R\&D Lead}
    {GenLots}
    {Lausanne, Switzerland}
    {September 2019 - March 2021}
    {
      \begin{cvitems}
        \item Built a microservices platform delivering client ERP data to production ML services, including reinforcement-learning-based supply-chain planning, \& led applied research with Swiss watch manufacturers
      \end{cvitems}
    }

\cventry
    {Research Scientist (PhD)}
    {University of Lausanne}
    {Lausanne, Switzerland}
    {November 2015 - August 2019}
    {
      \begin{cvitems}
        \item Led machine-learning \& privacy research on large-scale mobility data (300+ participant study) \& industry projects with public transport operators; 14 peer-reviewed publications
      \end{cvitems}
    }

\cventry
    {Project Engineer, ETH Zürich (2015){\enskip\cdotp\enskip}Embedded Systems Engineer, BOLT IoT (2011 - 2013)}
    {Earlier roles}
    {}
    {}
    {}

\end{cventries}
